\section{Background and related work}
\label{ch:background}

\subsection{Retrieval-augmented generation}

Retrieval-augmented generation conditions a language model's response on candidate evidence retrieved from an external collection. Lewis et al. formalize an influential architecture joining a learned retriever, a document index, and a generator~\cite{lewis2020rag}, providing a foundation for using external evidence in knowledge-intensive tasks. The chatbot uses this general principle but not the paper's joint-training procedure.

Three questions arise in such a system. Was useful evidence available in the collection? Did retrieval place it in the model's context? Did the answer use it correctly and cover the user's need? A model can fail despite receiving a relevant passage, and it can produce a plausible answer when retrieval supplies no adequate support, so retrieval results and answer quality are not interchangeable measures.

The distinction matters for procedural assistance. One passage may state a requirement while another describes an exception or where a form must be submitted; an answer that repeats only the first can be supported sentence by sentence yet remain incomplete. A question may also require information held in a private record, where public evidence explains the general process without determining the individual's outcome. RAG changes what information is available to generation; it does not grant institutional authority or access to records.

\subsection{Dense retrieval and candidate selection}

Dense retrieval represents questions and passages as learned numerical vectors. In the dual-encoder formulation described by Karpukhin et al., query and passage encoders operate separately and an inner product scores their compatibility~\cite{karpukhin2020dpr}. Passage vectors can be prepared in advance, so at request time only the question is encoded and nearby passages are retrieved as a limited candidate set.

For vectors produced by query encoder $f_q$ and passage encoder $f_d$, a simple retrieval score is
\[
  s_{\mathrm{dense}}(q,d)=f_q(q)^{\mathsf{T}}f_d(d).
\]
This expression illustrates the dual-encoder idea rather than every indexing choice in the implementation. Similarity measures, normalization, index configuration, and candidate limits also affect which passages reach the next stage, and must be recorded to reproduce a retrieval result.

Institutional abbreviations, dates, article numbers, and short named procedures may be important to the question. Lexical retrieval, in the tradition of the probabilistic relevance framework and BM25~\cite{robertson2009bm25}, can supply additional candidates when exact wording matters; the implementation studied here uses PostgreSQL full-text ranking for this role. Combining candidate sources requires deciding how to merge duplicates, distribute a limited candidate budget, and preserve source information before later scoring.

M3-Embedding addresses multilinguality, multiple retrieval functions, and different text lengths within one representation framework~\cite{chen2024m3}. Its dense, sparse, and multivector capabilities should be distinguished from the mode an application uses: the chatbot's evaluated path uses BGE-M3 query embeddings for its indexed evidence. The published capabilities motivate this choice of component, but the paper's benchmark results do not estimate retrieval effectiveness on FCSE regulations or student questions.

\subsection{Cross-encoder reranking and score thresholds}

A reranker revisits the retrieved candidates with a scoring function that inspects the question and passage together. Nogueira and Cho describe BERT passage reranking as a second stage after candidate retrieval~\cite{nogueira2019bertrerank}. The reranker reorders or filters an already limited pool; it cannot recover a passage absent from that pool.

The chatbot's retrieval path uses the multilingual BGE-reranker-v2-m3 model, whose official documentation describes query-passage scoring and optional sigmoid normalization~\cite{baaiBgeRerankerV2M3}. Experiment records, not the model card, identify the runtime artifacts that produced the saved scores.

A minimum score controls which candidates are admitted into answer context. If $r(q,d)$ is the reranker score and $c(d)$ denotes a source type, admission under an inclusive threshold convention can be written as
\[
  d\ \text{is eligible when}\ r(q,d)\geq\tau_{c(d)}.
\]
A shared threshold sets the same $\tau$ for every source type; separate thresholds permit different values for FAQs and document passages. Whether equality is admitted is an implementation choice that must remain fixed in a comparison, and a subsequent budget or ordering rule may still exclude an eligible candidate. The threshold is therefore one part of context selection, not a complete description of retrieval behavior.

Neither a raw score nor its sigmoid transform expresses the probability that a passage suffices to answer the question. Probability calibration and classification accuracy are distinct concepts~\cite{guo2017calibration}; placing a score between zero and one does not demonstrate calibrated answer support. In this thesis, thresholds are operational selection settings, not confidence percentages, and reranker probabilities are not calibrated.

\subsection{Heterogeneous FAQs and institutional documents}

An FAQ places a compact question near an answer, often in wording close to a common information need. A regulation follows institutional structure, with definitions, numbered provisions, qualifications, and amendments. The useful retrieval unit may therefore resemble a complete response in one corpus and only part of the necessary reasoning in the other, so the same selection rule need not have the same effect on both.

Document chunking makes longer sources searchable within a finite context budget but can separate a provision from its heading, exception, or continuation. Preserving section labels and retrieving neighboring chunks can restore local context, at the cost of enlarging the context with text that may contribute little. Chunk count and context length therefore describe evidence delivery rather than measure quality.

Authority and currentness add dimensions that similarity does not resolve. A passage may match the wording of a question while applying to another category of student or an earlier version of a procedure, and a correct source link does not prove that the selected rule applies. The architecture developed here therefore keeps source provenance throughout ingestion and retrieval so that answer support can be examined beyond a numerical ranking.

This heterogeneity motivates the source-threshold comparison. Lowering the document threshold may recover relevant passages whose scores do not resemble FAQ scores, or introduce ambiguity, repetition, and unrelated conditions. A best-candidate fallback raises a related choice when no evidence passes the threshold: an empty context becomes nonempty, but the best available candidate may still be unsuitable. Downstream assistance outcomes are needed to judge either policy.

\subsection{Context use and the limits of adding evidence}

Once material is selected, the generator must decide what bears on the question. Additional passages could supply a missing requirement, leave the answer unchanged, or blur the line between general and case-specific information.

Shi et al. show that irrelevant sentences can distract models solving arithmetic problems~\cite{shi2023distracted}. Hagstr\"om et al. study claim verification using realistically retrieved evidence and compare context utilization with synthetic settings~\cite{hagstrom2025reality}. Liu et al. find that models use information less reliably when it sits in the middle of a long input than near its beginning or end~\cite{liu2024lost}, so admitting more passages changes where key evidence appears as well as how much text surrounds it. These findings motivate checking actual evidence conditions rather than assuming that longer context helps or harms.

The present study therefore separates intervention from outcome. Context length, source counts, and the presence of reference quotations explain what the retrieval policy changed; supported answering, completeness, and appropriate routing describe the response. A match between a reference quotation and a retrieved block is useful trace evidence, not proof of semantic recall or correct application of the quoted provision.

\subsection{Macedonian and multilingual evaluation}

MIRACL provides a human-judged retrieval benchmark covering eighteen languages~\cite{zhang2023miracl}, but Macedonian is outside its coverage, and institutional documents and student questions differ from a benchmark collection in terminology, information needs, and source structure.

NoMIRACL examines multilingual relevance recognition using evidence pools that contain either no relevant passages or at least one relevant passage~\cite{thakur2024nomiracl}. This distinction is valuable for a system that must recognize inadequate evidence, although judging material claims in a generated answer is a different endpoint: a system can recognize relevant evidence and still omit a condition when answering.

Krsteski et al. develop Macedonian language-model and corpus resources and examine language-specific evaluation, including concerns about translated benchmarks~\cite{krsteski2025macedonian}. Their work supports examining the assistant in its own language and setting; handling Cyrillic text alone would not demonstrate command of institutional vocabulary or correct interpretation of a regulation.

FCSE student questions vary in script, transliteration, length, and detail. Script indicators describe visible text properties but do not identify language reliably: a Latin-script question could be English, transliterated Macedonian, or a mixture containing names and abbreviations. Differences between email and chatbot outcomes therefore cannot be attributed to language from script distributions alone. This thesis uses observed questions and limited script-related checks rather than a controlled language benchmark.

\subsection{Selective answering and useful guidance}

Kamath et al. study selective question answering, in which a system may abstain, under domain shift using a risk-coverage perspective for extractive question answering~\cite{kamath2020selective}. Coverage concerns how often the system answers, while risk concerns errors among the answers it gives. The work motivates evaluating abstention behavior.

Institutional assistance requires more than answer versus silence. A clarification can resolve an underspecified question, a referral can identify who can inspect a private record, and verification guidance can acknowledge uncertainty about a changing detail. These responses can be useful when no direct answer is justified, whereas refusing an adequately supported answer leaves an avoidable gap.

The assessment in this thesis therefore separates supported answering, useful safe routing, incompleteness, unnecessary abstention, unsupported claims, and uncertainty. A polite referral does not count as useful merely because it avoids a false statement: the next step must suit the information need, and the response must not imply that explaining a process has completed it.

\subsection{Automated assessment and confidence in labels}

Automated RAG evaluation can examine several dimensions of quality. RAGAs distinguishes faithfulness, answer relevance, and context relevance through automated measures~\cite{es2024ragas}. The present experiment uses its own frozen rubric rather than the RAGAs implementation.

ARES also separates context relevance, answer faithfulness, and answer relevance, and uses human-labeled validation data and prediction-powered inference~\cite{angelopoulos2023ppi} to address error in automated judgments~\cite{saadfalcon2024ares}. Resampling questions while retaining fixed AI labels, as done here, estimates variation across questions but does not provide this human-label correction.

Zheng et al. examine position and verbosity biases and failures in reasoning assessment, and distinguish consistency from accuracy~\cite{zheng2023judge}: a judge that repeats the same decision may consistently misread a source. Panickssery et al. find, on summarization tasks, that LLM evaluators can recognize their own outputs and tend to rate them more favorably~\cite{panickssery2024selfpreference}. Chen et al. further distinguish response quality from self-preference when models judge their own outputs~\cite{chen2025selfpreference}. Disagreement between two assessors can thus reflect capability, prompt interpretation, or reference quality as well as self-preference.

These issues apply directly when one model constructs references, generates answers, and supplies the primary assessment; the roles are separate in execution but may share errors. Mechanical checks can confirm that a quotation occurs in a named source or that a judgment follows the required schema, but not entailment, source applicability, or whether every necessary fact was identified. If a reference retains no facts, a completeness check becomes uninformative even though the record passes its structural checks. Repeated and cross-model assessment are therefore sensitivity analyses: they show whether categories and flags remain stable under another judgment, without creating an institutionally approved answer set.

\subsection{Reference architectures and RAG engineering}

A software reference architecture abstracts recurring responsibilities, components and interfaces from concrete systems in a domain. Angelov et al. analyze reference architectures along three dimensions, goal, context and design, and distinguish architectures whose goal is to standardize concrete architectures for interoperability from those that facilitate the design of new systems~\cite{angelov2012reference}. In their terms, the architecture in this thesis is a facilitation architecture for a single organization: it is derived from one implemented system and offered as a design aid, not as a standard validated across organizations. In design-science terms, it is an artifact whose utility must be demonstrated through evaluation~\cite{hevner2004design}; here that evaluation is the controlled study of one instantiation, together with an analysis of where its observed failures arise.

Engineering experience with RAG systems supports organizing such an architecture around failure-prone stages. Barnett et al. report seven failure points from three RAG case studies, including missing content, relevant documents ranked too low or excluded during context consolidation, answers that fail to extract information present in the context, and incomplete answers~\cite{barnett2024seven}. The positions of these failure points in their indexing and query pipeline correspond closely to the responsibilities separated in Chapter~\ref{ch:architecture}, which makes the observed failures in this study attributable to a stage rather than only to ``the chatbot''.

\subsection{University-support precedent and research position}

Marcel is a closely related conversational system for university student support~\cite{trienes2025marcel}. It combines retrieval from university resources with answer generation, evaluates anonymized student emails using manually prepared reference material, and reports deployment observations separately. It shows that real correspondence can support evaluation of university information needs; its manually prepared references contrast with the AI-created, unapproved references used here.

The present thesis does not introduce email-derived evaluation. Its contribution is the connection between an inspectable chatbot architecture, a controlled source-selection comparison, and an explicit analysis of assessment limitations. The cited work motivates each component, but results from other benchmarks, languages, and institutions do not transfer directly to FCSE; whether the selected evidence produces a complete answer or useful next step must be examined on the actual questions.
